% --- CORRECCIÓN IMPORTANTE EN LA PRIMERA LÍNEA ---
% Agregamos "xcolor=table" para que funcione \rowcolor sin errores
\documentclass[aspectratio=169, 10pt, xcolor=table]{beamer}

% --- Configuración del Tema (estilo profesional como el ejemplo) ---
\usetheme{Madrid}
\usecolortheme{dolphin}

% --- Paquetes necesarios ---
\usepackage[utf8]{inputenc}
\usepackage[spanish]{babel}
\usepackage{amsmath, amssymb}
\usepackage{booktabs}
\usepackage{graphicx}
\usepackage{listings}
\usepackage{tikz}
\usetikzlibrary{arrows.meta, positioning, shapes.geometric}

% --- Ruta de Imágenes (crea la carpeta "imagenes" y coloca ahí tus archivos) ---
\graphicspath{{./imagenes/}}

% --- Estilo para código Python ---
\definecolor{codegreen}{rgb}{0,0.6,0}
\definecolor{codegray}{rgb}{0.5,0.5,0.5}
\definecolor{codepurple}{rgb}{0.58,0,0.82}
\definecolor{backcolour}{rgb}{0.95,0.95,0.92}
\lstset{
    language=Python,
    backgroundcolor=\color{backcolour},
    commentstyle=\color{codegreen},
    keywordstyle=\color{blue},
    numberstyle=\tiny\color{codegray},
    stringstyle=\color{codepurple},
    basicstyle=\ttfamily\scriptsize,
    breaklines=true,
    frame=single,
    showstringspaces=false
}

% --- Pie de página personalizado ---
\makeatletter

\setbeamertemplate{footline}{
  \leavevmode%
  \hbox{%
  \begin{beamercolorbox}[wd=.333333\paperwidth,ht=2.25ex,dp=1ex,center]{author in head/foot}%
    \usebeamerfont{author in head/foot}\insertshortauthor
  \end{beamercolorbox}%
  \begin{beamercolorbox}[wd=.333333\paperwidth,ht=2.25ex,dp=1ex,center]{title in head/foot}%
    \usebeamerfont{title in head/foot}Sesión 8
  \end{beamercolorbox}%
  \begin{beamercolorbox}[wd=.333333\paperwidth,ht=2.25ex,dp=1ex,right]{date in head/foot}%
    \usebeamerfont{date in head/foot}Machine Learning \hfill \insertframenumber/\inserttotalframenumber \hspace*{0.3cm}
  \end{beamercolorbox}}%
  \vskip0pt%
}

\makeatother

% --- Datos de la portada ---
\title[SESIÓN 8]{\textbf{Sesión 8}}
\subtitle{Evaluación y Validación de Modelos: \\ Generalización, sesgo-varianza y optimización de hiperparámetros}
\author[Prof. C. Sánchez]{Carlos César Sánchez Coronel}
\institute{\textbf{MACHINE LEARNING}}
\date{}

% Logo opcional (descomentar si tienes la imagen)
% \titlegraphic{\includegraphics[height=1.5cm]{logo.png}}

% --- Eliminar la barra de navegación (opcional) ---
\setbeamertemplate{navigation symbols}{}

% --- Índice al inicio de cada sección ---
\AtBeginSection[]{
  \begin{frame}
    \frametitle{Contenido}
    \tableofcontents[currentsection]
  \end{frame}
}

% ============================================================================
% INICIO DEL DOCUMENTO
% ============================================================================
\begin{document}

% --- Portada ---
\begin{frame}
    \titlepage
\end{frame}

% --- Índice general ---
\begin{frame}{Contenido}
    \tableofcontents
\end{frame}

% ============================================================================
% SECCIÓN 1: UNDERFITTING Y OVERFITTING
% ============================================================================
\section{Underfitting y Overfitting}

\begin{frame}{Definición y causas}
    \begin{columns}[t]
        \column{0.5\textwidth}
        \begin{block}{Underfitting}
            \begin{itemize}
                \item Modelo demasiado simple para capturar la estructura.
                \item Alto error en entrenamiento y prueba.
                \item Causas: poca complejidad, características insuficientes, regularización excesiva.
            \end{itemize}
        \end{block}
        \column{0.5\textwidth}
        \begin{block}{Overfitting}
            \begin{itemize}
                \item Modelo se ajusta al ruido, no a la señal.
                \item Error muy bajo en entrenamiento, alto en prueba.
                \item Causas: modelo complejo, pocos datos, entrenamiento excesivo, falta de regularización.
            \end{itemize}
        \end{block}
    \end{columns}
    % Basado en Pregunta 1
    % IMAGEN: Gráfica de error vs complejidad
\end{frame}

\begin{frame}{Diagnóstico y soluciones}
    \begin{table}
        \centering
        \begin{tabular}{p{4cm}|p{6cm}}
            \toprule
            \textbf{Problema} & \textbf{Posibles soluciones} \\
            \midrule
            Underfitting & Aumentar complejidad, añadir características, reducir regularización, aumentar número de iteraciones. \\
            Overfitting & Reducir complejidad, aumentar regularización, más datos, early stopping, simplificar modelo. \\
            \bottomrule
        \end{tabular}
    \end{table}
    \begin{itemize}
        \item Herramientas gráficas: curvas de aprendizaje, curvas de validación.
    \end{itemize}
    % Basado en Pregunta 1
\end{frame}

% ============================================================================
% SECCIÓN 2: COMPROMISO SESGO-VARIANZA
% ============================================================================
\section{Compromiso Sesgo-Varianza}

\begin{frame}{Descomposición del error esperado}
    \begin{block}{Para un punto fijo $x_0$}
        \[
        \mathbb{E}\left[ (y_0 - \hat{f}(x_0))^2 \right] = \text{Sesgo}^2 + \text{Varianza} + \text{Ruido irreducible}
        \]
        \begin{itemize}
            \item \textbf{Sesgo}: error por suposiciones erróneas (diferencia entre predicción promedio y valor real).
            \item \textbf{Varianza}: variabilidad de las predicciones al cambiar el conjunto de entrenamiento.
            \item \textbf{Ruido}: error inherente a los datos, irreducible.
        \end{itemize}
    \end{block}
    % Basado en Pregunta 2
\end{frame}

\begin{frame}{Relación con la complejidad}
    \begin{itemize}
        \item Modelos simples $\rightarrow$ alto sesgo, baja varianza.
        \item Modelos complejos $\rightarrow$ bajo sesgo, alta varianza.
        \item El error total tiene un mínimo en el punto de equilibrio.
    \end{itemize}
    % IMAGEN: Gráfica típica de sesgo-varianza vs complejidad
    \vspace{0.5cm}
    \begin{exampleblock}{Diagnóstico práctico}
        \begin{itemize}
            \item Alto sesgo: errores altos en entrenamiento y validación.
            \item Alta varianza: error bajo en entrenamiento, alto en validación.
        \end{itemize}
    \end{exampleblock}
    % Basado en Pregunta 3
\end{frame}

% ============================================================================
% SECCIÓN 3: VALIDACIÓN CRUZADA
% ============================================================================
\section{Validación Cruzada (Cross-Validation)}

\begin{frame}{k-Fold Cross-Validation}
    \begin{columns}[t]
        \column{0.5\textwidth}
        \begin{block}{Procedimiento}
            \begin{enumerate}
                \item Dividir datos en $k$ pliegues.
                \item Para $i=1..k$: entrenar con $k-1$ pliegues, validar en el $i$-ésimo.
                \item Error promedio = $\frac{1}{k}\sum_{i=1}^k \text{error}_i$.
            \end{enumerate}
        \end{block}
        \column{0.5\textwidth}
        \begin{block}{Elección de $k$}
            \begin{itemize}
                \item $k=5$ o $10$ típicos.
                \item $k$ grande $\rightarrow$ menor sesgo, mayor varianza y costo.
                \item $k=n$ (LOOCV): máximo costo, alta varianza.
            \end{itemize}
        \end{block}
    \end{columns}
    % Basado en Pregunta 4
\end{frame}

\begin{frame}{Stratified k-Fold y LOOCV}
    \begin{itemize}
        \item \textbf{Stratified k-Fold}: mantiene proporción de clases en cada pliegue. Obligatorio en clasificación con desbalance.
        \item \textbf{Leave-One-Out (LOOCV)}: $k=n$. Entrena con todas menos una observación. Útil con pocos datos, pero costoso y de alta varianza.
    \end{itemize}
    % Basado en Preguntas 4 y 5
\end{frame}

\begin{frame}{Validación cruzada en series temporales}
    \begin{alertblock}{¡No usar k-fold estándar!}
        Viola la dependencia temporal (usa datos futuros para predecir pasado).
    \end{alertblock}
    \begin{columns}[t]
        \column{0.5\textwidth}
        \textbf{Expanding window}
        \begin{itemize}
            \item Entrenamiento crece, validación fija.
        \end{itemize}
        \column{0.5\textwidth}
        \textbf{Rolling window}
        \begin{itemize}
            \item Ventana de tamaño fijo se desliza.
        \end{itemize}
    \end{columns}
    \vspace{0.3cm}
    En scikit-learn: \texttt{TimeSeriesSplit}.
    % Basado en Pregunta 6
\end{frame}

\begin{frame}{Estimación del error y su varianza}
    \begin{exampleblock}{Ejemplo numérico}
        Errores en 5-fold: [0.15, 0.12, 0.18, 0.14, 0.16]
        \begin{itemize}
            \item Media = 0.15
            \item Desviación estándar $\approx$ 0.02
        \end{itemize}
        El error esperado es $0.15 \pm 0.02$.
    \end{exampleblock}
    % Basado en Pregunta 4
\end{frame}

% ============================================================================
% SECCIÓN 4: CURVAS DE APRENDIZAJE
% ============================================================================
\section{Curvas de Aprendizaje}

\begin{frame}{Construcción e interpretación}
    \begin{itemize}
        \item Muestran error en entrenamiento y validación vs. tamaño del conjunto de entrenamiento.
        \item Se construyen entrenando el modelo con subconjuntos crecientes.
    \end{itemize}
    \begin{columns}[t]
        \column{0.5\textwidth}
        \textbf{Alto sesgo}
        \begin{itemize}
            \item Curvas convergen a error alto.
            \item Añadir datos no ayuda.
        \end{itemize}
        \column{0.5\textwidth}
        \textbf{Alta varianza}
        \begin{itemize}
            \item Gran brecha entre curvas.
            \item Añadir datos reduce la brecha.
        \end{itemize}
    \end{columns}
    % Basado en Preguntas 3 y 7
    % IMAGEN: Curvas de aprendizaje típicas
\end{frame}

% ============================================================================
% SECCIÓN 5: BÚSQUEDA DE HIPERPARÁMETROS
% ============================================================================
\section{Búsqueda de Hiperparámetros}

\begin{frame}{Grid Search}
    \begin{block}{Definición}
        Prueba todas las combinaciones de una rejilla de valores.
    \end{block}
    \begin{itemize}
        \item Ventaja: exhaustivo.
        \item Desventaja: costo exponencial en número de parámetros.
    \end{itemize}
    \begin{exampleblock}{Ejemplo}
        \texttt{param\_grid = \{'max\_depth': [3,5,7], 'learning\_rate': [0.01,0.1,0.3], 'n\_estimators': [100,200]\}} \\
        $3\times3\times2 = 18$ combinaciones.
    \end{exampleblock}
    % Basado en Pregunta 8
\end{frame}

\begin{frame}{Random Search}
    \begin{itemize}
        \item Muestrea aleatoriamente combinaciones de una distribución.
        \item Más eficiente que Grid Search en espacios de alta dimensión.
        \item Con el mismo presupuesto, explora más valores por parámetro.
    \end{itemize}
    % Basado en Preguntas 8 y 9
\end{frame}

\begin{frame}{Bayesian Optimization}
    \begin{block}{Idea}
        Construye un modelo probabilístico (Gaussian Process) de la función objetivo y lo usa para decidir el próximo punto a evaluar.
    \end{block}
    \begin{itemize}
        \item Función de adquisición: \textit{Expected Improvement}, \textit{Upper Confidence Bound}.
        \item Requiere menos evaluaciones que Random Search, ideal para modelos costosos.
    \end{itemize}
    Implementaciones: Optuna, Hyperopt, Scikit-Optimize.
    % Basado en Pregunta 10
\end{frame}

\begin{frame}{Comparación de métodos}
    \begin{table}
        \centering
        \begin{tabular}{l|c|c|c}
            \toprule
            \textbf{Método} & \textbf{Exhaustividad} & \textbf{Eficiencia} & \textbf{Aplicación típica} \\
            \midrule
            Grid Search & Alta & Baja & Pocos parámetros, valores discretos \\
            Random Search & Media & Media & Muchos parámetros, presupuesto limitado \\
            Bayesian Opt. & Adaptativa & Alta & Modelos costosos, muchas evaluaciones \\
            \bottomrule
        \end{tabular}
    \end{table}
    % Basado en teoría y Pregunta 19
\end{frame}

% ============================================================================
% SECCIÓN 6: CONSIDERACIONES PRÁCTICAS
% ============================================================================
\section{Consideraciones Prácticas}

\begin{frame}{Data Leakage (Fuga de datos)}
    \begin{alertblock}{Error común}
        Usar información del conjunto de prueba o validación durante el entrenamiento.
    \end{alertblock}
    \begin{exampleblock}{Ejemplos}
        \begin{itemize}
            \item Estandarizar con media y desviación de todo el dataset antes de dividir.
            \item Hacer oversampling (SMOTE) antes de dividir.
            \item Seleccionar características usando todo el dataset.
            \item Target encoding con media global.
        \end{itemize}
    \end{exampleblock}
    % Basado en Preguntas 11 y 18
\end{frame}

\begin{frame}{Cómo evitar leakage}
    \begin{enumerate}
        \item Siempre dividir primero: entrenamiento / validación / prueba.
        \item Ajustar transformaciones (escalado, imputación, codificación) solo en entrenamiento, luego aplicar a validación/prueba.
        \item Usar validación cruzada para selección de características o target encoding.
        \item En series temporales, respetar el orden.
        \item En datos con grupos, usar \texttt{GroupKFold}.
    \end{enumerate}
    % Basado en Preguntas 11 y 12
\end{frame}

\begin{frame}{Validación con grupos y entrenamiento final}
    \begin{itemize}
        \item \textbf{GroupKFold}: asegura que todas las muestras de un mismo grupo (ej. paciente) estén en el mismo pliegue.
        \item Una vez elegidos hiperparámetros, entrenar el modelo final con \textbf{todos los datos disponibles} (entrenamiento + validación) y evaluar una sola vez en el conjunto de prueba.
    \end{itemize}
    % Basado en Preguntas 12 y 13
\end{frame}

% ============================================================================
% SECCIÓN 7: CASO INTEGRADO: XGBOOST
% ============================================================================
\section{Caso Integrado: XGBoost}

\begin{frame}{Selección de hiperparámetros para XGBoost}
    \begin{itemize}
        \item Problema: clasificación binaria, 50,000 muestras, 20 características, clases balanceadas.
        \item Pasos:
        \begin{enumerate}
            \item Dividir en entrenamiento (70\%), validación (15\%), prueba (15\%).
            \item Elegir hiperparámetros a optimizar: \texttt{learning\_rate}, \texttt{max\_depth}, \texttt{subsample}, \texttt{colsample\_bytree}, \texttt{lambda}, \texttt{alpha}.
            \item Usar \texttt{RandomizedSearchCV} o \texttt{Optuna} con validación cruzada (5-fold).
            \item Evaluar el mejor modelo en el conjunto de prueba.
        \end{enumerate}
    \end{itemize}
    % Basado en teoría y Pregunta 20
\end{frame}

% ============================================================================
% SECCIÓN 8: MÉTRICAS EN VALIDACIÓN CRUZADA
% ============================================================================
\section{Métricas en Validación Cruzada}

\begin{frame}{Métricas según el problema}
    \begin{itemize}
        \item \textbf{Clasificación}: accuracy, AUC, F1, log-loss. Elegir según desbalance y objetivos.
        \item \textbf{Regresión}: RMSE, MAE, RMSLE, R².
        \item En validación cruzada, se promedia la métrica y se reporta su desviación.
    \end{itemize}
    \begin{exampleblock}{Ejemplo con alta varianza}
        Modelo A: errores [0.10,0.12,0.11,0.13,0.09] (media=0.11, sd=0.016) \\
        Modelo B: errores [0.05,0.20,0.04,0.22,0.06] (media=0.114, sd=0.083) \\
        Se prefiere el modelo A por estabilidad.
    \end{exampleblock}
    % Basado en Pregunta 14
\end{frame}

% ============================================================================
% SECCIÓN 9: CONEXIÓN CON EL RESTO DEL CURSO
% ============================================================================
\section{Conexión con el resto del curso}

\begin{frame}{Importancia de la validación}
    \begin{itemize}
        \item Todas las sesiones anteriores (regresión, clasificación, árboles, boosting) requieren evaluación rigurosa.
        \item La elección de hiperparámetros es crítica para el rendimiento final.
        \item Técnicas vistas aquí se aplican a cualquier modelo supervisado o no supervisado.
        \item Conceptos de sesgo-varianza subyacen a la regularización y complejidad.
    \end{itemize}
\end{frame}

% ============================================================================
% SECCIÓN 10: RESUMEN
% ============================================================================
\section{Resumen}

\begin{frame}{Lo que hemos aprendido}
    \begin{exampleblock}{Conceptos clave}
        \begin{itemize}
            \item \textbf{Underfitting / Overfitting}: diagnóstico y soluciones.
            \item \textbf{Sesgo-Varianza}: descomposición del error y trade-off.
            \item \textbf{Validación cruzada}: k-fold, stratified, LOOCV, time series.
            \item \textbf{Curvas de aprendizaje}: identifican sesgo/varianza.
            \item \textbf{Búsqueda de hiperparámetros}: Grid, Random, Bayesian.
            \item \textbf{Data leakage}: cómo evitarlo.
            \item \textbf{Grupos y evaluación final}.
        \end{itemize}
    \end{exampleblock}
\end{frame}

% --- Diapositiva final ---
\begin{frame}
    \centering
    \Huge \textbf{¡Gracias!}

\end{frame}

\end{document}